\begin{frame}{이번 과목의 정체성}
  \begin{block}{한 줄씩}
    \kb{강화학습} = 행동하고, 보상만 보고 배우는 학습. \\[0.6em]
    \kb{이번 학기의 무대} = \textbf{실제 응용}(로봇팔 모방학습, 자율주행 등)을
    \textbf{물리 시뮬레이션}으로 재현한 환경.
  \end{block}
  \vskip1em
  ``학습을 배우는 학습''이라는 큰 그림을 먼저 머릿속에 걸어두면,
  후반부의 DQN$\cdot$PPO$\cdot$MCTS가 별개의 알고리즘이 아니라
  \textbf{같은 틀의 다른 적용}으로 읽힘.
  \begin{itemize}
    \item 지도의 정답 $y$가 \kb{보상} $R$으로 바뀜.
    \item 비지도의 ``구조''가 \kb{에피소드}(행동-보상 시퀀스)로 바뀜.
    \item 경사하강법은 \textbf{그대로} --- 다만 경사가
      \textit{에이전트의 경험}에서 옴.
  \end{itemize}
\end{frame}
